% ABOUTME: Evaluation protocol section of the IEEE report (issue #2).
% ABOUTME: Engine level splitting, leakage, metric definitions and the early vs late stage split.
\section{Evaluation Protocol}
\label{sec:evaluation}

\subsection{Splitting by engine}
Each training engine in FD001 contributes one row per operating cycle, and
consecutive cycles of the same engine differ only by a small amount of
degradation and sensor noise. If rows were split at random, cycle $t$ of an
engine could fall in the training set while cycle $t+1$ fell in the
validation set. The model would then be scored on near copies of rows it had
already seen, and the validation error would understate the error on a new
engine. We therefore split by engine: every row of a given \texttt{unit} lies
on one side of the split only. Hyperparameters are tuned with five fold
\texttt{GroupKFold} cross validation grouped on \texttt{unit}, with a fixed
seed of 42. The test set is used once, after all hyperparameters are frozen.

\subsection{Scoring on test engines}
Training engines run to failure, but each test engine is cut off at an
unknown point before failure, and the dataset gives the true RUL only at its
last observed cycle. Each test engine is therefore scored once, using the
prediction at its final cycle, so all 100 test engines carry equal weight.

\subsection{Metrics}
Let $y_i$ be the true RUL of engine $i$, $\hat{y}_i$ the prediction,
$d_i = \hat{y}_i - y_i$ the error and $n$ the number of engines. We report
\begin{align}
\mathrm{RMSE} &= \sqrt{\tfrac{1}{n}\textstyle\sum_{i} d_i^2}, \qquad
\mathrm{MAE} = \tfrac{1}{n}\textstyle\sum_{i} \lvert d_i \rvert, \\
S &= \sum_{i} s_i, \quad
s_i =
\begin{cases}
e^{-d_i/13} - 1, & d_i < 0,\\
e^{d_i/10} - 1, & d_i \ge 0,
\end{cases}
\end{align}
where $S$ is the asymmetric scoring function of the PHM08 challenge
\cite{saxena2008damage}. RMSE penalises large errors more than MAE does,
and both treat early and late predictions alike. The score $S$ does not: a
prediction that overestimates RUL by 10 cycles costs as much as one that
underestimates it by 13, and the penalty grows exponentially. Overestimating
RUL means the engine is believed healthy when it is close to failure, which
risks an in service failure; underestimating it only brings maintenance
forward. Because $S$ is a sum rather than a mean, it is compared only across
models scored on the same set of engines.

\subsection{Early and late stage error}
Maintenance decisions are made near the end of an engine's life, so error
there matters more than error far from failure. We split the test engines by
true RUL into a late stage ($y_i \le 50$ cycles) and an early stage
($y_i > 50$) and report RMSE and MAE for each. A model with a good overall
RMSE but a poor late stage RMSE would be of little use for scheduling
maintenance.
